\documentclass[11pt, a4paper]{article}

% Gemeinsame Style-Definitionen (TEXINPUTS muss templates/ enthalten — siehe scripts/build_tex.py)
\input{bfw-style.tex}

\title{\vspace*{-1.5cm}\Huge\bfseries\color{bfwPrimaryDark}Session~6: Wirtschaftskreislauf \& Marktsituationen\\[0.4em]
  \large\color{bfwInkSoft}\textnormal{Geldströme, Marktformen und Preisbildung im Marktumfeld von JIKU IT-Solutions}}
\author{\textsf{IT Lernfeld 1 --- Das Unternehmen und die eigene Rolle im Betrieb beschreiben}}
\date{}

\begin{document}
\maketitle
\thispagestyle{empty}

\vspace{1em}
\begin{merksatz}[title={\bfseries Worum geht es in dieser Session?}]
Unternehmen agieren nicht im Vakuum --- sie sind in wirtschaftliche Kreisläufe eingebunden,
in denen Geld, Güter und Faktorleistungen zwischen Haushalten, Unternehmen, Banken, Staat
und Ausland fließen. Am Beispiel von \emph{JIKU IT-Solutions} analysieren wir, in welchen
Märkten das Unternehmen aktiv ist, welche Marktform vorherrscht, wie sich Preise durch
Angebot und Nachfrage bilden und welche Rahmenbedingungen das Marktverhalten beeinflussen.
\end{merksatz}

\tableofcontents
\vspace{1em}
\hrule
\vspace{1em}

\section{Wirtschaftskreisläufe}

\subsection{Einfacher Wirtschaftskreislauf}

Der einfache Wirtschaftskreislauf beschränkt sich auf zwei Wirtschaftssubjekte: \textbf{Haushalte}
und \textbf{Unternehmen}. Haushalte stellen den Unternehmen volkswirtschaftliche
Faktorleistungen --- Arbeit, Boden und Kapital --- zur Verfügung, damit diese Konsumgüter
erstellen können. Als Gegenleistung zahlen die Unternehmen Löhne, Gehälter, Miete, Pacht
oder Zinsen an die Haushalte. Die Haushalte nutzen dieses Einkommen, um Konsumgüter zu
kaufen, woraus die Konsumausgaben entstehen.

\begin{merksatz}[title={\bfseries Merke: Gegenläufige Ströme}]
Im Wirtschaftskreislauf fließen \textbf{Geld- und Güterströme stets in entgegengesetzten
Richtungen}. Wo Güter oder Faktorleistungen hinfließen, fließt Geld in die entgegengesetzte
Richtung.
\end{merksatz}

\begin{center}
\begin{tabular}{p{3.5cm} p{9.5cm}}
\toprule
\textbf{Strom} & \textbf{Richtung und Inhalt}\\
\midrule
Faktorleistungen & Haushalte → Unternehmen: Arbeit, Boden, Kapital\\
\addlinespace
Einkommen & Unternehmen → Haushalte: Löhne, Gehälter, Zinsen, Miete, Pacht\\
\addlinespace
Konsumgüter & Unternehmen → Haushalte: Waren und Dienstleistungen\\
\addlinespace
Konsumausgaben & Haushalte → Unternehmen: Zahlungen für Konsumgüter\\
\bottomrule
\end{tabular}
\end{center}

\subsection{Erweiterter Wirtschaftskreislauf}

Die volkswirtschaftliche Realität ist wesentlich komplexer. Der erweiterte Wirtschaftskreislauf
bezieht zusätzlich \textbf{Staat}, \textbf{Banken} und \textbf{Ausland} als Wirtschaftssubjekte ein.

\subsubsection*{Rolle der Banken}

Haushalte legen Ersparnisse bei Banken an und erhalten dafür Zinsen. Banken nutzen diese
Mittel, um Unternehmen Investitionskredite (sog.\ Deckungskredite) zu gewähren. Auch der
Staat kann sich über Banken und den Geldmarkt verschulden.

\subsubsection*{Rolle des Staates}

\begin{center}
\begin{tabular}{p{4cm} p{8.5cm}}
\toprule
\textbf{Geldstrom} & \textbf{Erläuterung}\\
\midrule
Direkte Steuern & Haushalte → Staat: z.\,B. Kfz-Steuer, Grundsteuer, Einkommensteuer\newline
{\footnotesize (Lohnsteuer wird vom Arbeitgeber einbehalten und abgeführt — am Abführungsweg ändert das nichts an der Einordnung als direkte Steuer)}\\
\addlinespace
Indirekte Steuern & Staat vereinnahmt über Unternehmen: z.\,B. Umsatzsteuer, Mineralölsteuer\\
\addlinespace
Transferleistungen & Staat → Haushalte: z.\,B. Kindergeld, Sozialhilfe, Sparzulage\\
\addlinespace
Staatsausgaben & Staat → Unternehmen/Infrastruktur: Bildung, Verkehr, Sicherheit\\
\addlinespace
Subventionen & Staat → Unternehmen: Investitionsförderung, Arbeitsplatzschaffung\\
\bottomrule
\end{tabular}
\end{center}

\begin{merksatz}[title={\bfseries Merke: Direkte vs. indirekte Steuern}]
\textbf{Direkte Steuern} entrichten Haushalte unmittelbar an den Staat (z.\,B. Kfz-Steuer,
Grundsteuer). \textbf{Indirekte Steuern} zahlen Verbraucher zunächst an Unternehmen, die
sie dann an den Staat abführen (z.\,B. Umsatzsteuer, Mineralölsteuer). Beide Steuerarten
erhöhen die Staatsquote.
\end{merksatz}

\subsubsection*{Rolle des Auslands}

Im Zuge der Globalisierung nimmt das Ausland eine wichtige Rolle ein. Deutschland ist
einerseits stark exportorientiert --- Exporterlöse fließen aus dem Ausland nach Deutschland.
Andererseits bedeuten Importausgaben und Kapitalanlagen im Ausland einen Abfluss von Geld.
Durch das weltweite Bankennetz sind heute internationale Kapitaltransfers in Sekunden möglich.

\section{Marktsituationen beschreiben}

\subsection{Marktarten und Transaktionsbereiche}

Unternehmen wie JIKU IT-Solutions müssen entscheiden, in welchen Märkten sie aktiv sein wollen.
Märkte lassen sich nach verschiedenen Kriterien unterscheiden:

\begin{center}
\begin{tabular}{p{4cm} p{5cm} p{4cm}}
\toprule
\textbf{Marktart} & \textbf{Erläuterung} & \textbf{Beispiel}\\
\midrule
Freier Markt & Jeder hat Zugang zum Markt & Automarkt, Lebensmittelmarkt\\
\addlinespace
Geschlossener Markt & Nicht jeder hat Zugang & Waffenmarkt, Großhandelsmarkt\\
\addlinespace
Gütermarkt & Verbrauch/Gebrauch durch Endverbraucher oder Investoren & IT-Markt, Maschinenmarkt\\
\addlinespace
Faktorenmarkt & Handel mit Produktionsfaktoren & Arbeitsmarkt, Finanzmärkte\\
\addlinespace
Vollkommener Markt & Alle Bedingungen erfüllt & Buttermarkt, Milchmarkt\\
\addlinespace
Unvollkommener Markt & Bedingungen nicht erfüllt & Benzinmarkt, Schmuckmarkt\\
\bottomrule
\end{tabular}
\end{center}

Im ITK-Bereich sind Transaktionsbereiche wichtig. Je nach Vertragspartner unterscheidet man:

\begin{center}
\begin{tabular}{p{2cm} p{11cm}}
\toprule
\textbf{Kürzel} & \textbf{Bedeutung und Beispiel}\\
\midrule
B2B & Business-to-Business: Unternehmen handelt mit Unternehmen, z.\,B. JIKU liefert Server an Firmenkunden\\
B2C & Business-to-Consumer: Unternehmen bietet Verbrauchern Produkte an, z.\,B. Online-Shop\\
C2C & Consumer-to-Consumer: Verbraucher verkauft an Verbraucher, z.\,B. Gebrauchtteile-Plattform\\
A2B / A2C & Administration-to-Business/Consumer: Behörden digitalisieren Leistungen (E-Government)\\
\bottomrule
\end{tabular}
\end{center}

\subsection{Bedingungen des vollkommenen Marktes}

\begin{center}
\begin{tabular}{p{4cm} p{9cm}}
\toprule
\textbf{Bedingung} & \textbf{Erläuterung}\\
\midrule
Vollkommene Konkurrenz & Viele Anbieter und Nachfrager treffen sich auf dem Markt\\
\addlinespace
Güterhomogenität & Es werden nur gleichartige Güter angeboten\\
\addlinespace
Markttransparenz & Marktteilnehmer haben den Überblick über alle Anbieter und Preise\\
\addlinespace
Präferenzlosigkeit & Nachfrager haben keine besonderen Vorlieben für bestimmte Anbieter\\
\addlinespace
Anpassungsfähigkeit & Marktteilnehmer passen ihr Verhalten sofort dem Preis an\\
\bottomrule
\end{tabular}
\end{center}

\section{Marktformen}

Marktformen klassifizieren Märkte nach der \textbf{Anzahl der Anbieter und Nachfrager} und
dem daraus resultierenden Einfluss auf die Preisbildung.

\begin{center}
\begin{tabular}{p{2cm} p{3.5cm} p{3.5cm} p{3.5cm}}
\toprule
\textbf{} & \textbf{Einer (Anbieter)} & \textbf{Wenige (Anbieter)} & \textbf{Viele (Anbieter)}\\
\midrule
\textbf{Einer (Nachfr.)} & Zweiseitiges Monopol & Beschränktes Nachfragemonopol & Nachfragemonopol\\
\addlinespace
\textbf{Wenige (Nachfr.)} & Beschränktes Angebotsmonopol & Zweiseitiges Oligopol & Nachfrageoligopol\\
\addlinespace
\textbf{Viele (Nachfr.)} & Angebotsmonopol & Angebotsoligopol & Polypol\\
\bottomrule
\end{tabular}
\end{center}

\begin{merksatz}[title={\bfseries Merke: Lesart der Matrix}]
Die Zeilen zeigen die \textbf{Nachfrageseite}, die Spalten die \textbf{Angebotsseite}.
Im Feld oben links: ein Anbieter und ein Nachfrager = zweiseitiges Monopol. Unten rechts:
viele Anbieter und viele Nachfrager = Polypol.
\end{merksatz}

\subsection{Monopol}

Ein Monopolist ist der einzige Anbieter (Angebotsmonopol) oder der einzige Nachfrager
(Nachfragemonopol) auf einem Markt. Er kann seine Marktposition ausnutzen und den Preis
bestimmen. Unternehmen versuchen auch in polypolistischen Märkten, monopolähnliche
Positionen zu erlangen (\emph{Quasimonopol}), z.\,B. durch Patente, Netzwerkeffekte,
Diversifikation oder Exklusivität.

Nach \textbf{§\,18 Abs.\,4 GWB} (Gesetz gegen Wettbewerbsbeschränkungen) wird eine
marktbeherrschende Stellung vermutet, wenn ein Unternehmen mindestens 40\,\% Marktanteil
hält. Das Bundeskartellamt prüft dann, ob Maßnahmen gegen Missbrauch zu ergreifen sind.

\begin{merksatz}[title={\bfseries § 19 GWB --- Missbrauchsverbot}]
Die missbräuchliche Ausnutzung einer marktbeherrschenden Stellung durch ein oder mehrere
Unternehmen ist verboten. Ein Unternehmen ist marktbeherrschend, soweit es als Anbieter
oder Nachfrager ohne Wettbewerb oder ohne wesentlichen Wettbewerb ist.
\end{merksatz}

\subsection{Oligopol}

Das \textbf{Angebotsoligopol} ist eine Marktform, bei der wenigen Anbietern viele Nachfrager
gegenüberstehen. Die Anbieter beobachten sich gegenseitig und reagieren aufeinander
(Parallelverhalten). Beispiele im ITK-Bereich: Chipmarkt (Intel, AMD, Qualcomm), Cloud-Anbieter
(AWS, Azure, Google Cloud), Smartphone-Betriebssysteme (Android/iOS), Spielkonsolenmarkt.

\subsection{Polypol}

Das \textbf{Polypol} ist die vom Staat bevorzugte Marktform: Einer Vielzahl von Anbietern
steht eine Vielzahl von Nachfragern gegenüber. Durch das Konkurrenzverhalten bildet sich
der Marktpreis durch freies Angebot und freie Nachfrage. Beispiele: Lebensmittelmärkte,
App-Entwickler-Markt, IT-Freelancer.

\section{Preisbildung durch Angebot und Nachfrage}

\subsection{Gleichgewichtspreis}

Im vollkommenen Markt orientieren sich alle Marktteilnehmer am Preis. Die
\textbf{Nachfragekurve} fällt mit steigendem Preis (bei höheren Preisen kaufen Kunden
weniger). Die \textbf{Angebotskurve} steigt mit dem Preis (bei höheren Preisen lohnt sich
mehr Produktion). Im \textbf{Schnittpunkt} beider Kurven liegt der Gleichgewichtspreis und
die Gleichgewichtsmenge.

Beispiel aus dem Lehrbuch: Bei einem Preis von 450,00\,€ sind angebotene und nachgefragte
Menge mit je 650 Stück identisch --- das ist der Gleichgewichtspreis.

\subsection{Angebotsüberhang und Nachfrageüberhang}

\begin{center}
\begin{tabular}{p{4.5cm} p{4cm} p{4.5cm}}
\toprule
\textbf{Situation} & \textbf{Verhältnis} & \textbf{Folge}\\
\midrule
Angebotsüberhang & Angebot > Nachfrage & Käufermarkt: Preisdruck nach unten\\
\addlinespace
Nachfrageüberhang & Nachfrage > Angebot & Verkäufermarkt: Preisdruck nach oben\\
\addlinespace
Gleichgewicht & Angebot = Nachfrage & Stabiler Gleichgewichtspreis\\
\bottomrule
\end{tabular}
\end{center}

\subsection{Kurvenverschiebungen}

Externe Faktoren können die Angebots- oder Nachfragekurve verschieben und so einen neuen
Gleichgewichtspreis erzeugen:

\begin{itemize}
  \item \textbf{Nachfragekurve verschiebt sich nach rechts} (N0 → N1): z.\,B. durch
        Einkommenserhöhung oder Steuersenkung → Gleichgewichtspreis steigt.
  \item \textbf{Angebotskurve verschiebt sich nach links} (A0 → A1): z.\,B. durch höhere
        Produktionskosten oder schlechte Konjunkturerwartungen → Gleichgewichtspreis steigt.
\end{itemize}

\section{Marktwirtschaftliche Rahmenbedingungen}

Die Marktbedingungen werden durch drei Arten von Veränderungen geprägt:

\begin{center}
\begin{tabular}{p{3cm} p{5cm} p{5cm}}
\toprule
\textbf{Art} & \textbf{Ursachen} & \textbf{Beispiele}\\
\midrule
Saisonal & Wetter, Jahreszeiten, Kalender & Weihnachtsgeschäft, Ferienzeiten, Fastenzeiten\\
\addlinespace
Strukturell & Neue Technologien, Gesetze, Wertvorstellungen & IT-Innovationen, Umweltgesetze, Trend zu Nachhaltigkeit\\
\addlinespace
Konjunkturell & BIP-Schwankungen im Wirtschaftszyklus & Aufschwung → Boom → Rezession → Depression\\
\bottomrule
\end{tabular}
\end{center}

\begin{merksatz}[title={\bfseries Zusammenfassung: Kernbegriffe dieser Session}]
\textbf{Wirtschaftskreislauf}: Geld- und Güterströme fließen gegenläufig zwischen Haushalten,
Unternehmen, Staat, Banken und Ausland.~·~
\textbf{Marktform}: Klassifikation nach Anbieter- und Nachfragerzahl (Monopol, Oligopol,
Polypol).~·~
\textbf{§\,18 Abs.\,4 GWB}: ab 40\,\% Marktanteil gilt marktbeherrschende Stellung.~·~
\textbf{Gleichgewichtspreis}: Schnittpunkt von Angebots- und Nachfragekurve.~·~
\textbf{Käufermarkt}: Angebot > Nachfrage; \textbf{Verkäufermarkt}: Nachfrage > Angebot.
\end{merksatz}

\end{document}
